\documentclass{book}
\makeatletter
\@ifundefined{frontmatter}{\newcommand{\frontmatter}{\clearpage}}{}
\@ifundefined{mainmatter}{\newcommand{\mainmatter}{\clearpage}}{}
\@ifundefined{backmatter}{\newcommand{\backmatter}{\clearpage}}{}
\makeatother
\title{Book Full Fixture}
\author{Ferritex Bench}
\begin{document}
\frontmatter
\maketitle
\tableofcontents

\chapter{Preface}
The front matter provides a short preface that introduces the corpus document. It is intentionally brief but includes inline math $u_n = u_{n-1} + 1$ to mix narrative and notation.

\mainmatter
\chapter{Foundations}\label{chap:foundations}
This chapter introduces the main discussion and refers ahead to Chapter \ref{chap:applications}.

\section{Concepts}
We use the classical identity
\begin{equation}
\sum_{k=1}^{n} k = \frac{n(n+1)}{2}
\end{equation}
to keep the content recognizable while exercising numbered displays.

\section{Context}
The prose is intentionally plain: the goal is to combine book structure, chapter navigation, and equation layout in one self-contained source.

\chapter{Applications}\label{chap:applications}
This chapter points back to Chapter \ref{chap:foundations} and adds a second layer of sectioning.

\section{Case Study}
Assume that $p(x) = x^3 - x$. Then
\[
p(2) = 6
\]
and the text around the formula verifies that ordinary paragraphs and displayed math coexist correctly.

\section{Observations}
The book fixture uses front matter, main matter, back matter, a table of contents, chapters, sections, and math in a single source file.

\backmatter
\chapter{Closing Notes}
Back matter closes the document with a concise summary of the integration features covered by the fixture.
\end{document}
